(1) 单位长度KVL给 \(\partial_xV=-R'I-L'\partial_tI\)，KCL给 \(\partial_xI=-G'V-C'\partial_tV\)，由此可求二阶电报方程 \(V_{xx}=L'C'V_{tt}+(R'C'+L'G')V_t+R'G'V\)，电流有相同形式。

(2) 对 \(e^{i\omega t-\gamma x}\) 相量，\(\gamma^2=(R'+i\omega L')(G'+i\omega C')\)，\(\boxed{Z_0=\sqrt{(R'+i\omega L')/(G'+i\omega C')}}\)。对沿正方向传播的被动线选 \(\operatorname{Re}\gamma\ge0\)，\(Z_0\) 的分支要求实功率前向流动。

(3) 若 \(R'/L'=G'/C'=a\)，则 \(R'+i\omega L'=L'(a+i\omega)\)、\(G'+i\omega C'=C'(a+i\omega)\)，故 \(\gamma=\sqrt{L'C'}(a+i\omega)\)，相位随频率严格线性，幅度衰减 \(e^{-a\sqrt{L'C'}x}\) 不依赖频率。\(Z_0=\sqrt{L'/C'}\) 是实常数，群速与相速均为 \(1/\sqrt{L'C'}\)。

(4) 无损极限 \(R'=G'=0\) 时 \(\gamma=i\omega\sqrt{L'C'}\)、\(Z_0=\sqrt{L'/C'}\)。若电感可忽略且 \(G'=0\)，则 \(\gamma\simeq\sqrt{i\omega R'C'}\)，回到扩散型RC传输线。热耗散为每长度 \(R'|I|^2/2+G'|V|^2/2\ge0\)；两项分别反映金属与绝缘损耗。
